\section{ALEPH detector}
\label{sec:aleph}

ALEPH was a general-purpose LEP detector with near $4\pi$ coverage and high-precision charged-particle tracking in a 1.5~T solenoidal magnetic field~\cite{Decamp:1990jra,ALEPH:1994ayc}. The detector subsystems most relevant for this event-shape analysis are the central tracker, electromagnetic and hadronic calorimeters, and outer muon chambers.

The central tracking system combined a silicon vertex detector, an inner tracking chamber, and a large time projection chamber (TPC). Together these detectors provided precise track parameters and momentum measurements, and they define the object-quality observables used in this analysis (for example $d_0$, $z_0$, and $N_{\mathrm{TPC}}$). The charged-track momentum resolution is one of the key ingredients controlling event-shape resolution at detector level.

The calorimeter system comprised two main subsystems. These are the electromagnetic calorimeter (ECAL) and the hadronic calorimeter (HCAL). ECAL information dominates the response to photons and electromagnetic showers. HCAL is essential for neutral-hadron energy flow and global event-energy reconstruction. The muon system outside HCAL supports lepton categorization in the archived ALEPH object definitions.

In this analysis, we use the archived ALEPH reconstructed object representation directly, including per-particle category labels (\texttt{pwflag}) and kinematic variables. Event-shape observables are built from selected charged and neutral particles, so detector performance in each reconstructed particle category is a direct input to thrust migration and unfolding stability.
